% © 2026 Ing. David Jaroš — CC BY-NC-ND 4.0
% Canonical covariant-tetrad reframe, July 2026.
% UBT-AI-PROVENANCE-BEGIN
% schema: ubt-ai-provenance/v1
% tier: B_machine_verified
% ai_assistance: disclosed
% human_review: machine-verification
% editorial_responsibility: Ing. David Jaroš
% policy: ../../AI_PROVENANCE.md
% notice: Machine-verified against named sources or verifiers; individual attestation is not claimed.
% UBT-AI-PROVENANCE-END

\documentclass[12pt]{article}
\usepackage[a4paper,margin=2.5cm]{geometry}
\usepackage{amsmath,amssymb,amsthm,mathtools}
\usepackage[most]{tcolorbox}
\usepackage[colorlinks=true,linkcolor=blue]{hyperref}
\usepackage{ubtprovenance}
\UBTTier{B}
\newtheorem{definition}{Definition}[section]
\newtheorem{theorem}[definition]{Theorem}
\newtheorem{proposition}[definition]{Proposition}
\newtheorem{corollary}[definition]{Corollary}
\theoremstyle{remark}\newtheorem{remark}[definition]{Remark}
\title{Covariant Biquaternionic Tetrads in UBT:\\Central Anticommutator Metric and the Local Rank Theorem}
\author{Ing. David Jaroš}
\date{15 July 2026}
\begin{document}\maketitle
\UBTProvenanceNotice
\begin{abstract}
We formulate the local classical geometry of Unified Biquaternion Theory without
an embedding map, a compact-fiber average, or an externally chosen projection.
The covariant first jet $E_\mu:=\mathcal N_0^{-1/2}D_\mu\Theta$ is restricted in
the classical sector to a real Lorentz slice of the biquaternions.  On this
slice the symmetrised product with quaternion conjugation is automatically
central and equals a Lorentzian scalar times the algebra unit.  This defines
the metric intrinsically.  We prove that the tetrad-to-metric differential has
rank ten at every nondegenerate tetrad, with a six-dimensional kernel equal to
local Lorentz frame rotations.  The old comparison ``eight real components of
$\Theta$ versus ten metric components'' therefore is not a valid local
kinematic obstruction: the relevant object is the first covariant jet.  The
remaining problem is dynamical and integrability-based: derive the connection,
preserve the Lorentz slice, and solve $D_\mu\Theta=E_\mu$ on shell.
\end{abstract}

\section{Algebraic setting}
Let
\[
 \mathbb B=\mathbb C\otimes\mathbb H
\]
with commuting complex unit $i$ and quaternion units
$\mathbf e_j\mathbf e_k=-\delta_{jk}\mathbf1+\epsilon_{jk\ell}\mathbf e_\ell$.
Denote quaternion conjugation, which reverses the three quaternion units but
does not conjugate the commuting complex coefficient, by $\sharp$:
\[
 (z^0\mathbf1+z^k\mathbf e_k)^\sharp
 =z^0\mathbf1-z^k\mathbf e_k.
\]

\begin{definition}[Lorentz slice]
The classical Lorentz module is the real four-dimensional subspace
\[
 W_L:=\{\,i x^0\mathbf1+x^k\mathbf e_k\mid x^a\in\mathbb R\,\}
 \subset\mathbb B.
\]
For $X=i x^0\mathbf1+x^k\mathbf e_k$ and
$Y=i y^0\mathbf1+y^k\mathbf e_k$, define
\[
 \eta(X,Y):=-x^0y^0+\sum_{k=1}^3x^ky^k.
\]
\end{definition}

\begin{theorem}[Central Jordan identity]
For all $X,Y\in W_L$,
\begin{equation}
 \boxed{\frac12\left(X^\sharp Y+Y^\sharp X\right)
       =\eta(X,Y)\mathbf1.}
 \label{eq:central-jordan}
\end{equation}
Thus the symmetrised product is already a real central biquaternion; no trace,
real-part map, or fiber projection is required.
\end{theorem}

\begin{proof}
Write $X=(ix^0,\mathbf x)$ and $Y=(iy^0,\mathbf y)$ in scalar--vector
quaternion notation, for which
$(a,\mathbf u)(b,\mathbf v)=(ab-\mathbf u\cdot\mathbf v,
 a\mathbf v+b\mathbf u+\mathbf u\times\mathbf v)$.
Then $X^\sharp=(ix^0,-\mathbf x)$.  The scalar part of $X^\sharp Y$ is
$-x^0y^0+\mathbf x\cdot\mathbf y$.  Its vector part is
$i x^0\mathbf y-i y^0\mathbf x-\mathbf x\times\mathbf y$.
The vector part of $Y^\sharp X$ is its negative, while the scalar part is the
same.  Their half-sum is therefore exactly \eqref{eq:central-jordan}.
\end{proof}

\section{Covariant first jet and metric}
The fundamental UBT field remains $\Theta(q,\tau)\in\mathbb B$.  On physical
spacetime it is a smooth field.  Let $D_\mu$ be the covariant derivative in the
representation carried by $\Theta$ and let $\mathcal N_0>0$ be a fixed global
unit-setting constant.

\begin{definition}[UBT tetrad sector]
Define
\begin{equation}
 E_\mu:=\mathcal N_0^{-1/2}D_\mu\Theta.
 \label{eq:E-def}
\end{equation}
The classical tetrad sector consists of configurations for which
$E_\mu\in W_L$ and the four $E_\mu$ are linearly independent over $\mathbb R$.
The metric is the unique real coefficient in the central identity
\begin{equation}
 \boxed{\frac12\left(E_\mu^\sharp E_\nu+E_\nu^\sharp E_\mu\right)
       =g_{\mu\nu}\mathbf1.}
 \label{eq:metric-central}
\end{equation}
\end{definition}

Writing $E_\mu=i e_\mu{}^0\mathbf1+e_\mu{}^k\mathbf e_k$ gives immediately
\begin{equation}
 g_{\mu\nu}=e_\mu{}^a e_\nu{}^b\eta_{ab},
 \qquad \eta_{ab}=\operatorname{diag}(-1,1,1,1).
 \label{eq:tetrad-metric}
\end{equation}
No local normalization is permitted; such a denominator would alter the
field equations and can freeze a metric component.  $\mathcal N_0$ is constant
and may be absorbed into units.

\begin{definition}[Algebraic bivector companion]
The antisymmetric companion of the metric is
\begin{equation}
 \Sigma_{\mu\nu}:=\frac12\left(E_\mu^\sharp E_\nu-E_\nu^\sharp E_\mu\right),
 \qquad \Sigma_{\nu\mu}=-\Sigma_{\mu\nu}.
 \label{eq:bivector}
\end{equation}
It contains oriented-plane/spin information.  It must not be confused with
the curvature $[D_\mu,D_\nu]$.
\end{definition}

\section{Local rank theorem}
Let $e=(e_\mu{}^a)$ be a nondegenerate tetrad and define
$\Phi(e)=e\eta e^{\mathsf T}$.

\begin{theorem}[Surjectivity of the tetrad-to-metric differential]
At every invertible real tetrad, the differential
\[
 d\Phi_e:\mathbb R^{4\times4}\longrightarrow\operatorname{Sym}_4(\mathbb R)
\]
has rank ten.  Its kernel has dimension six and consists precisely of
infinitesimal local Lorentz rotations of the tetrad.
\end{theorem}

\begin{proof}
The first variation is
\begin{equation}
 \delta g_{\mu\nu}=\eta_{ab}
 \left(\delta e_\mu{}^a e_\nu{}^b+e_\mu{}^a\delta e_\nu{}^b\right).
 \label{eq:metric-variation}
\end{equation}
Given an arbitrary symmetric tensor $h_{\mu\nu}$, choose
\begin{equation}
 \delta e_\mu{}^a=\frac12 h_{\mu\rho}e^{\rho a},
 \label{eq:surjective-choice}
\end{equation}
where $e^{\rho a}$ is the inverse tetrad with its Lorentz index raised.  Direct
substitution into \eqref{eq:metric-variation} gives
$\delta g_{\mu\nu}=h_{\mu\nu}$.  Hence $d\Phi_e$ is surjective and has rank ten.
The domain has dimension sixteen, so the kernel has dimension six.  Variations
$\delta e=e\lambda$ with $\lambda^{\mathsf T}\eta+\eta\lambda=0$ lie in the
kernel and form the six-dimensional Lorentz Lie algebra; therefore they exhaust
it.
\end{proof}

\begin{corollary}[Resolution of the naive component-count obstruction]
The local algebraic map $E_\mu\mapsto g_{\mu\nu}$ has full metric rank.  The
comparison $\dim_\mathbb R\mathbb B=8<10$ is irrelevant to this map because the
metric is determined by four covariant first derivatives, not by the value of
$\Theta$ alone.  This closes only the \emph{kinematic tetrad rank} question.
\end{corollary}

\section{Connection, Christoffel symbols, and flat limit}
A connection is the rule needed to compare internal frames at neighboring
points.  In representation-independent notation,
\begin{equation}
 D_\mu\Theta=\partial_\mu\Theta+\rho_*(\Omega_\mu)\Theta,
 \label{eq:D-general}
\end{equation}
where $\Omega_\mu$ is a biquaternionic/Lorentz-frame connection and $\rho_*$
specifies how it acts on $\Theta$.  This notation deliberately does not choose
left, right, or adjoint action before the representation is fixed.

The Christoffel symbols $\Gamma^\rho{}_{\mu\nu}$ and $\Omega_\mu$ are related
but are not the same object.  $\Gamma$ transports coordinate indices; $\Omega$
transports the local biquaternionic/Lorentz frame.  Their compatibility is the
tetrad postulate
\begin{equation}
 \boxed{\partial_\mu E_\nu-\Gamma^\rho{}_{\mu\nu}E_\rho
       +\rho_{\mathrm{vec}*}(\Omega_\mu)E_\nu=0.}
 \label{eq:tetrad-postulate}
\end{equation}
For a chosen torsion condition and a nondegenerate tetrad, this equation makes
$\Gamma$ and $\Omega$ two descriptions of the same local geometry.

The connection curvature is
\begin{equation}
 \mathcal R_{\mu\nu}(\Omega)
 :=[D_\mu,D_\nu]
 =\partial_\mu\Omega_\nu-\partial_\nu\Omega_\mu
  +[\Omega_\mu,\Omega_\nu]
 \label{eq:connection-curvature}
\end{equation}
with the last expression understood in the selected representation.
In a flat simply connected region, $\mathcal R_{\mu\nu}=0$ permits a local
frame in which $\Omega_\mu=0$.  In inertial Cartesian coordinates one also has
$\Gamma^\rho{}_{\mu\nu}=0$, and $D_\mu=\partial_\mu$: this is the special
relativistic sector.  In curvilinear coordinates $\Gamma$ and $\Omega$ may be
nonzero even when the curvature vanishes.

\section{The remaining nontrivial bridge}
The tetrad is not postulated independently; UBT requires
$E_\mu=D_\mu\Theta/\sqrt{\mathcal N_0}$.  Therefore full closure needs more
than the rank theorem.  One must prove:
\begin{enumerate}
 \item existence of a connection and field solving $D_\mu\Theta=\sqrt{\mathcal N_0}E_\mu$;
 \item preservation of the Lorentz slice $E_\mu\in W_L$ and centrality of
       \eqref{eq:metric-central};
 \item derivation or unique elimination of $\Omega_\mu$ from the UBT action,
       so it is not an arbitrary extra field;
 \item compatibility/integrability.  In a torsion-free coordinate sector,
       antisymmetrisation gives schematically
       \[
       D_\mu E_\nu-D_\nu E_\mu=[D_\mu,D_\nu]\Theta
       =\mathcal R_{\mu\nu}(\Omega)\Theta;
       \]
 \item derivation of Einstein dynamics (or a definite UBT modification) from
       the same action.
\end{enumerate}
These are dynamical questions, not a shortage of local metric components.

\begin{tcolorbox}[colback=yellow!8!white,colframe=orange!70!black,
 title=Proof status]
\textbf{Proved:} central Lorentzian anticommutator identity; local tetrad
representation of every Lorentz metric; rank ten of the tetrad-to-metric map;
flat inertial limit $\Omega=\Gamma=0$ when curvature is zero.

\textbf{Open:} derivation of the admissible Lorentz slice and the connection
from the single fundamental field/action; global integrability; on-shell
representation of all required GR solutions; Einstein dynamics from the
canonical UBT action.

\textbf{Noncanonical alternative:} compact-$\psi$ fiber averaging remains a
mathematically valid candidate completion but is not used in this theorem and
is not part of the minimal covariant-tetrad core.
\end{tcolorbox}

\end{document}
